\begin{table}[h]
    \caption{Adopted churn definitions and associated studies.}
    \label{tab:RS_RQ1_churn_definition}
    \centering
    %{\fontfamily{phv}\selectfont   % inizio gruppo → Helvetica/Arial-like
    \scriptsize                 % o \footnotesize, \small…
    \resizebox{0.7\textwidth}{!}{
    \begin{tabular}{|
  >{\centering\arraybackslash}m{4.5cm}
  |>{\centering\arraybackslash}m{5cm}|}
    \hline
    \rowcolor[HTML]{D5E8D4} 
    \rule[-1.5ex]{0pt}{4ex}\textbf{Churn definition}
  & \rule[-1.5ex]{0pt}{4ex}\textbf{Study ID(s)} \\ \hline
    Static & S3, S5, S10, S13, S14, S15, S20, S21, S22, S25, S27, S39, S41 \\ \hline
    Dynamic & S6, S19, S30, S31, S33, S35, S36, S37 \\ \hline
    Static with dynamic window & S16, S23 \\ \hline
    Dynamic with dynamic window & S24 \\ \hline
    Cluster/Segment based & S9 \\ \hline
    Probabilistic based & S8 \\ \hline
    \end{tabular}
    }
 % }% fine gruppo → torna al font normale qui
\end{table}